\subsection{Graph comparisons for larger opposite blocks}

\begin{lemma}[Proved: graph realization with collapsed edges]
\label{CP2:A-directed-graph-collapses-eight}
In Lemma~\ref{CP2:A-directed-graph-right-zero-realization}, allow a
right monoid actor to send an oriented edge either to an oriented edge
with the same parameter, or to a vertex, in which case both endpoints
map to that vertex.  The natural equivariant graph realization remains
valid.  Consequently, if the acting monoid has order at most eight and
$Y=\operatorname{hocofib}(\Gamma_{hM}\to BM)$ has bounded rational
homology, every rational degree-two product on $Y$ vanishes.
\end{lemma}
\begin{proof}
Use exactly the state set consisting of vertices and edges, with the
two right-zero actors sending an edge to its endpoints and fixing
vertices.  A collapsed edge is sent to its image vertex as a state.
The endpoint equations say precisely that this set map commutes with
both right-zero actors.  On the graph category it sends the edge
object and both endpoint arrows to the image vertex and its identity.
Thus it defines an actual graph-category functor.  On the action
category, the image endpoint arrows are the corresponding right-zero
endomorphisms of that vertex; the graph functor sends them to the
identity.  Hence the graph comparison still commutes strictly with
these actors.  Its comma categories and their contractions are
unchanged.  This proves the asserted equivariant equivalence, including
loops, parallel edges and collapsed edges.

The resulting actual collage has opposite corner $R_3$.  The proof of
Theorem~\ref{CP2:A-directed-graph-eight} classifies the possible minimal
bands of $M$ and applies the all-state $R_3$ theorems to this collage;
none of those theorems requires its right state maps to preserve the
set of nonvertex states.  The same proof therefore gives the conclusion.
\end{proof}

\begin{proposition}[Proved: an order-eight corner opposite a graph block]
\label{CP2:A-semisimple-height-one-block-eight}
Let $C$ be a one-way collage of a corner monoid $M$, with $|M|\leq8$,
and a nonempty full block $D$.  Suppose all rational corner algebras
of $D$ are semisimple and its strict support has height at most one.
If $BC$ has bounded rational homology, all rational degree-two products
on $BC$ vanish.  Empty connecting components are permitted, provided
the whole connecting profunctor is nonempty.
\end{proposition}
\begin{proof}
The proof of Lemma~\ref{CP2:A-baseline-semisimple-pair-graph} applies
objectwise to this entire block.  There is one vertex set of orbits
for each object of $D$ and one balanced edge set for each nonempty
strict hom-set of $D$.  Each edge joins its two prescribed orbit
vertices.  There are no two composable strict arrows.  Quotienting
corner arrows and the balanced relations therefore gives an actual
functor from the category of elements to this graph.  Filtering by
strict object changes gives only columns zero and one.  Semisimplicity
makes every internal derived tensor concentrate in degree zero, so
the actual functor is a rational homology equivalence.  It is natural
for maps of acts; in particular the commuting $M$ action gives strict
maps of the graph preserving endpoint types.  This also proves
$H_j(BD;\mathbb Q)=0$ for $j\geq2$, even if $D$ is disconnected.

Collapse $BD$ in the actual two-block model and call the cofiber $Y$.
The quotient induces a surjection $H^2(Y;\mathbb Q)\to H^2(BC;\mathbb Q)$
and passes bounded rational homology to $Y$.  The equivariant graph
comparison remains a rational homology equivalence after taking monoid
homotopy orbits, by its nonnegative free bar filtration.  Its induced
map of actual cofibers therefore identifies the rational cohomology
ring of $Y$ with the graph cofiber.  The preceding lemma annihilates
its degree-two products.  Pullback proves the assertion for $BC$.
Taking opposites proves the other direction of the one-way collage.
\end{proof}

\begin{proposition}[Proved: a star with an order-eight center and small leaves]
\label{CP2:A-star-eight-three-leaves}
Let $C$ have source-star or sink-star support, with center corner of
order at most eight and every leaf corner of order at most three.
For arbitrary nonempty stalks, bounded rational homology of $BC$
forces every rational degree-two product to vanish.
\end{proposition}
\begin{proof}
By opposition take the center $M$ to be a source.  Every leaf corner
$N$ is rationally acyclic.  If its trivial right rational module is
projective, its left-action category maps naturally and rationally
equivalently to the discrete orbit set.  This is the actual orbit
comparison of Lemma~\ref{CP2:A-high-orbits}.  If that module is not
projective, the minimal-ideal projectivity test in order at most three
leaves precisely $N=R_3$: a nonprojective right augmentation needs at
least two minimal columns, and their band plus a separate identity
already uses all three elements.  Its action category has the natural
actual graph model of Lemma~\ref{CP2:A-high-P5-graph}.  A commuting
actor sends an edge either to the edge of its image state or to a
vertex.  The preceding lemma allows both possibilities.

The full leaf block is a disjoint union.  Its bridge is naturally a
disjoint union of the preceding action-category pieces, and hence has
an equivariant rational comparison to a finite graph.  The leaf block
has zero positive rational homology.  Collapsing it therefore passes
boundedness and surjects on second cohomology, just as in the preceding
proposition.  Apply Lemma~\ref{CP2:A-directed-graph-collapses-eight}
and pull back the zero products.  This proves the source-star case,
and opposition proves the sink-star case.
\end{proof}

\begin{corollary}[Proved: two gluing consequences]
\label{CP2:A-nineteen-graph-gluing-gates}
Each of the following conditions on a category with acyclic support
and bounded rational homology implies zero rational degree-two products.
\begin{enumerate}
\item The undirected comparability graph is a tree; every internal
corner has order at most five; and every edge satisfies one of the
bounded cup-zero two-object conditions in
Lemma~\ref{CP2:A-five18-tree-edges}, or has corner orders at most eight
with at least one of them at most three.
\item An object $p$ is strictly adjacent only to $q$, the corner at
$p$ has order at most eleven, the corner at $q$ has order at most two,
and all corners of the full deletion of $p$ are semisimple over
$\mathbb Q$ with strict support height at most three.
\end{enumerate}
In any of these assertions, one may first repeatedly delete singleton
pendants whose neighbors have corner order at most four.
\end{corollary}
\begin{proof}
For the first assertion, the new two-object edge condition is the
preceding star proposition with one leaf.  The actual tree-edge nerve
union and its corner intersections are the same as in
Lemma~\ref{CP2:A-five18-tree-edges}.  Corners of order at most five have
zero rational homology in degrees at least three, so boundedness
passes to every edge piece and degree-four restriction is injective.
The edge products vanish, proving the claim.

For the second assertion, write $E$ for the deletion and $F$ for the
full edge on $p,q$.  The nerves are the actual union $NE\cup NF$
with intersection the corner at $q$.  That corner is rationally
acyclic.  Thus boundedness passes to $BF$ and restriction in degree
four is injective.  The semisimple-height bound gives $H^4(BE;\mathbb Q)=0$.
At the endpoint $q$, the actual projective orbit functor identifies
the rational cohomology ring of $BF$ with that of a bounded ordinary
act cone over the corner at $p$, or its opposite.  The latter corner
has order at most eleven; Theorem~\ref{CP2:A-bounded-eleven-acts}
annihilates its degree-two products.  Hence both restrictions of any
product are zero.

Finally, each stated singleton-pendant deletion induces a rational
homology equivalence and rational cohomology-ring isomorphism by
Lemma~\ref{CP2:A-rational-pendant-deletion}.  It preserves boundedness
and transfers the established zero products back to the original
category.  This last step uses no integral or homotopy-equivalence
claim for that deletion.
\end{proof}
